
\subsection{Results and discussion}

\paragraph{Prongs close at the first attempt.} The prong set flow reproduces every prong-level marginal to
$W_1 \le 0.006$ (direction, both momentum hypotheses, proton score), the fractions of prongs with
kinematics, with a proton fit and exiting to within 0.3\%, and the per-event counts of each category
(Fig.~\ref{fig:m3-event}). The closure classifier on event-level prong summaries reaches AUC 0.516 without
and 0.511 with the truth features (Table~\ref{tab:m3-closure}), i.e.\ the prong content is
indistinguishable from MasterAnaDev at the level these summaries probe. The masked set flow with
cross-attention to the particle tokens therefore handles the variable-length output as designed in
Section~\ref{sec:m2}; no ordering, matching or slot assignment was needed.

\paragraph{The vertex comb, in three steps.} The discrete vertex head went through the same closure-driven
iteration as M2. With classes ``unsnapped'' and ``nearest plane $+\delta$, $|\delta|\le 3$'' the snapped fraction
came out at 55\% against 69\% in MasterAnaDev, because 14\% of reconstructed vertices snap to a plane more
than three planes away; a ninth class, decoded as the plane nearest to the flow's continuous $z$, fixed the
fraction (0.689 vs 0.690) and the offset-to-plane distribution ($W_1 = 0.07$\,mm,
Fig.~\ref{fig:m3-vertex}). The truth-conditional test on the vertex $z$ stayed at 0.578 until the head was given the
\emph{phase} of the true $z$ within the 22\,mm plane cycle as an explicit feature; the encoder cannot resolve a
22\,mm periodicity from a coordinate standardised over 4600\,mm. With it the test drops to 0.547 and the
head's log loss falls below the marginal (1.573 vs 1.621, evaluated with sampled flags and
multiplicity as conditions).

\paragraph{Effect on Tier 1.} Treating the snapped $z$ as a dequantised target for the flow removed the comb
from the continuous model. The Tier 1 closure improved from 0.57 / 0.62 (M2) to
0.534 / 0.534 (reco vector / truth $+$ reco), with every block at or below 0.525
conditionally; the M2 exit criterion of 0.55 is now met in both forms. The Tier 0 heads and the multiplicity
head are unchanged by the joint fine-tuning.

\paragraph{What is not modelled.} No prong-to-true-particle correspondence is produced: the flow learns the
prong set given the particle set, and any assignment lives implicitly in the cross-attention. The tuple's
own truth match (PDG, energy fraction) is kept in the training data for validation only. A per-prong pointer
head over the particle tokens would provide it and is the natural next addition if an analysis needs
truth-matched prongs from the surrogate. The two calibrated recoil variants remain derived quantities.

\paragraph{Deliverable.} \texttt{scripts/surrogate\_to\_ntuple.py} now writes the full pruned tuple (80 branches):
Tier 0 flags, muon, plane-snapped vertex, calorimetry, multiplicity, and the \texttt{pion\_*[10]},
\texttt{proton\_*} and \texttt{sec\_protons\_*} blocks through the canonical prong table of Section~\ref{sec:setup},
with the primary proton chosen as the prong with the highest generated primary score among those with a proton
fit.
